%% kapa: exam=1 round=23 year=2012 period=2 form=A subject=appraisal_law subject_ko=부동산관계법규 seq=2 number=42
%% source: Q-Net 감정평가사 기출문제 (cst00309) artlSeq=5052580 file=raw/1008/5052580/2012년 제23회 감정평가사 제1차시험 기출문제(2교시 A형).pdf
%% conversion: kordoc 4.24.51 -> tools/gen1.py (auto); review=proofread-sonnet; flags=none
\begin{question}{42}
국토의 계획 및 이용에 관한 법령상 개발행위허가의 절차 등에 관한 설명으로 옳지 않은 것은?
\begin{choices}
  \item 도시계획사업에 의한 개발행위는 허가를 받을 필요가 없다.
  \item 허가권자가 개발행위허가에 조건을 붙이려는 때에는 미리 신청한 자의 의견을 들어야 한다.
  \item 허가권자는 개발행위허가의 신청에 대하여 특별한 사유가 없으면 15일(심의 또는 협의기간은 제외) 이내에 처분을 하여야 한다.
  \item 허가권자가 처분을 할 때에는 허가증을 발급하거나 불허가처분의 사유를 서면으로 신청인에게 알려야 한다.
  \item 개발밀도관리구역 안에서 개발행위허가 신청을 할 때에는 기반시설의 설치나 그에 필요한 용지의 확보에 관한 계획서를 제출하여야 한다.
\end{choices}
\end{question}
\answer{5}
